\documentclass[10pt,letterpaper]{article}
\usepackage[margin=1in]{geometry}
\usepackage{times}
\usepackage{enumitem}
\usepackage{amsmath,amssymb}
\usepackage{xcolor}
\usepackage{hyperref}

\newcommand{\eg}{e.g.}
\begin{document}

\section*{3.1\quad Finding 1: Cross-Indicator Inconsistency}

A common assumption in existing pruning methods is that a single importance
indicator suffices to identify the most valuable visual tokens.
We challenge this assumption by systematically comparing indicators across
four representative methods that span three distinct paradigms:
\textit{vision-intrinsic saliency} (VisionZip~[54], using ViT CLS attention),
\textit{text-conditioned hybrid} (CDPruner~[58], coupling text--visual cosine
similarity with a conditional DPP kernel),
\textit{pure semantic diversity} (DivPrune~[3], using max-min pairwise distance),
and \textit{saliency--coverage hybrid} (SCOPE~[11], fusing CLS attention with
submodular coverage gain).

As visualized in Fig.~3, these indicators produce strikingly different token
selections---even on the same image under the same compression ratio.
VisionZip concentrates on regions receiving high CLS attention, typically
semantically salient foreground objects, \emph{irrespective of the user's
question}.
CDPruner shifts the selection toward regions most relevant to the text query,
while its DPP kernel simultaneously promotes geometric spread among the
retained tokens.
DivPrune, by contrast, distributes selections to maximize the minimum pairwise
distance in the feature space, yielding broad spatial coverage at the cost of
foreground detail.
SCOPE occupies a middle ground: its submodular coverage term prevents token
clustering in a narrow semantic region, while the CLS saliency weight steers
selection toward informative patches.

This divergence arises because the four methods optimize along fundamentally
\textit{orthogonal indicator dimensions}:

\begin{itemize}[leftmargin=*]
  \item \textbf{Dimension~A\,---\,Saliency} evaluates individual token
        importance and further decomposes into two sub-dimensions with
        distinct behaviors:
  \begin{itemize}
    \item \textbf{A1 (Vision-Intrinsic):}
          \textit{``Is this token intrinsically important in the image,
          regardless of the question?''}---relying solely on vision-encoder
          signals ({\eg}, CLS attention in VisionZip and SCOPE), producing
          a fixed ranking for a given image.
    \item \textbf{A2 (Task Relevance):}
          \textit{``Is this token relevant to the user's current
          query?''}---requiring text--visual interaction ({\eg}, CDPruner
          computes cosine similarity between each visual token and the
          instruction embedding), yielding query-dependent rankings that
          differ across prompts.
  \end{itemize}
  \item \textbf{Dimension~B\,---\,Representativeness} evaluates whether the
        selected subset collectively represents the original token set, and
        further decomposes along two complementary spaces:
  \begin{itemize}
    \item \textbf{B1 (Semantic Diversity):}
          \textit{``Are the selected tokens semantically non-redundant in the
          embedding space?''}---enforced via set-theoretic objectives that
          penalize feature-space duplication ({\eg}, DivPrune's MMDP maximizes
          the minimum pairwise distance; SCOPE's submodular coverage maximizes
          the sum of nearest-neighbor similarities; CDPruner's conditional DPP
          maximizes the volume spanned by the selected feature vectors).
    \item \textbf{B2 (Spatial Coverage):}
          \textit{``Does the selected set preserve at least one representative
          token from every spatial region of the image?''}---enforced via
          geometric partitioning of the token grid ({\eg}, Nuwa~[39] divides
          the grid into $M{\times}M$ non-overlapping blocks and selects top-$n$
          tokens per block; HoloV~[44] allocates per-crop budgets proportional
          to holistic scores).
          Notably, B1 and B2 are not redundant: a semantically diverse subset
          may still concentrate in a single spatial quadrant, while a spatially
          uniform selection may retain redundant tokens from texturally
          homogeneous regions.
  \end{itemize}
\end{itemize}

Critically, high saliency does not imply high representativeness: a set of
tokens may all receive top CLS attention scores yet cluster within a narrow
semantic region, leaving background context---essential for holistic
reasoning---entirely unrepresented.
SCOPE's own analysis confirms this: under pure saliency selection, the
$\theta$-coverage of the retained set can fall \emph{below random
sampling}~[11].
Conversely, maximizing diversity alone (DivPrune) scatters tokens across
task-irrelevant regions while omitting fine-grained details critical for
query-specific understanding.
The B2 dimension introduces yet another failure mode: methods operating purely
in the embedding space (B1) can produce spatially skewed selections that miss
entire image regions---a systematic weakness for OCR, spatial reasoning, and
referring expression comprehension tasks~[39].
These observations motivate a unified framework that jointly optimizes across
all four sub-dimensions.

\end{document}
